% ============================================================================
%  Example paper -- packages, marking roles, and settings
%
%  Every package below ships with the TeX Live collections that CI installs
%  (base, latex-recommended, latex-extra, plain-generic, pictures,
%  bibtex-extra, science).
%  Check the collection before adding a package.
% ============================================================================

% ---------- Page layout ----------
\usepackage[a4paper,margin=0.85in]{geometry}

% ---------- Typography ----------
\usepackage{microtype}
\usepackage{csquotes}

% ---------- Math ----------
\usepackage{amsmath,amssymb,amsthm,mathtools,bm}

% ---------- Graphics & figures ----------
\usepackage{graphicx}
\usepackage{tikz}
\usetikzlibrary{positioning,calc,arrows.meta}
\usepackage[font=small,labelfont=bf]{caption}
\usepackage{subcaption}

% ---------- Tables ----------
\usepackage{booktabs}
\usepackage{array}
\usepackage{tabularx}

% ---------- Lists ----------
\usepackage{enumitem}
\setlist{nosep, leftmargin=*}

% ---------- SI units ----------
\usepackage{siunitx}

% ---------- Marking module dependencies ----------
\usepackage[dvipsnames]{xcolor}
\usepackage{soul}
\usepackage{environ}
\usepackage{etoolbox}
\usepackage{xparse}
\usepackage[most]{tcolorbox}
\usepackage{xpatch}
% PDF annotations for \MarkEdit popup notes. pdfcomment loads hyperref itself,
% so it comes first to avoid a hyperref option clash.
\usepackage{pdfcomment}

% ---------- Hyperlinks & cross-references ----------
\usepackage{hyperref}
\hypersetup{
  colorlinks=true,
  linkcolor=blue,
  citecolor=ForestGreen,
  urlcolor=blue,
  breaklinks=true
}
\usepackage[nameinlink]{cleveref}
\crefname{equation}{Eq.}{Eqs.}
\Crefname{equation}{Eq.}{Eqs.}
\crefname{figure}{Fig.}{Figs.}
\Crefname{figure}{Fig.}{Figs.}

% ---------- Bibliography (biblatex + biber) ----------
% datamodel points at markings.dbx, which declares the markroles field.
\usepackage[
  backend=biber,
  style=numeric-comp,
  maxcitenames=2,
  maxbibnames=99,
  giveninits=true,
  doi=true,
  isbn=false,
  url=true,
  sorting=none,
  datamodel=../modules/markings/markings
]{biblatex}
\addbibresource{references.bib}

% ==== Marking module ====
% If you move this paper, update both this path and the datamodel path above.
% ============================================================================
% Marking module
% Public API:
%   \DefineMarkRole{role}{color}[state]   state: colored (default) | plain | hidden
%   \SetMarkFallbackColor{color}
%   \Mark{role}[kind]{content}
%   \SetMarkEditMode{plain|popup|inline}   default: plain
%   \MarkEdit{role}[kind]{old}{new}        kind: text (default) | math
%   \begin{MarkEnv}{role}[kind] ... \end{MarkEnv}
%   \MarkPlaceholderFigure{role}[layout]{description}{label}{caption}
% Supported Mark kinds: text, math, figure, heading
%
% Citation highlighting (biblatex only):
%   \SetMarkDefaultCiteColor{color}   (optional override -- rarely needed)
%   In .bib: markroles = {role}       (drives all highlighting automatically)
% ============================================================================

% Fail fast when required dependencies are not loaded in config.
\makeatletter
\@ifpackageloaded{xcolor}{}{
  \PackageError{marking}{MarkErrMissingDependency: xcolor not loaded}{}
}
\@ifpackageloaded{soul}{}{
  \PackageError{marking}{MarkErrMissingDependency: soul not loaded}{}
}
\@ifpackageloaded{environ}{}{
  \PackageError{marking}{MarkErrMissingDependency: environ not loaded}{}
}
\@ifpackageloaded{etoolbox}{}{
  \PackageError{marking}{MarkErrMissingDependency: etoolbox not loaded}{}
}
\@ifpackageloaded{xparse}{}{
  \PackageError{marking}{MarkErrMissingDependency: xparse not loaded}{}
}
\@ifpackageloaded{tcolorbox}{}{
  \PackageError{marking}{MarkErrMissingDependency: tcolorbox not loaded}{}
}
\makeatother

% Popup edit notes are PDF annotations. The mode remains plain unless selected.
\RequirePackage{pdfcomment}
% soul needs to treat citations as one unit in edit highlights and strikeouts.
\soulregister\cite7

\providecommand{\MarkingResolvedColor}{red}

\makeatletter
\newcommand{\marking@iftypesetTF}[2]{%
  \ifx\protect\@typeset@protect
    #1%
  \else
    #2%
  \fi
}
\makeatother

\ExplSyntaxOn
% Declare the expl3 argument variants this module uses. Recent expl3 releases
% predeclare them, so omitting these lines happens to work on an up-to-date
% TeX Live; older bundles (and Tectonic's pinned bundle) raise "Undefined
% control sequence" at the first \DefineMarkRole. Generating a variant that
% already exists is a no-op, so this is safe on every release.
\cs_generate_variant:Nn \tl_if_eq:nnTF { Vn }
\cs_generate_variant:Nn \tl_if_eq:nnT  { Vn }
\cs_generate_variant:Nn \tl_if_eq:nnF  { Vn }
\cs_generate_variant:Nn \tl_if_blank:nTF { V }
\cs_generate_variant:Nn \str_case:nnF { Vn }

\prop_new:N \g_marking_role_color_prop
\prop_new:N \g_marking_role_state_prop
\tl_new:N \g_marking_fallback_color_tl
\tl_gset:Nn \g_marking_fallback_color_tl { red }
\tl_new:N \g_marking_edit_mode_tl
\tl_gset:Nn \g_marking_edit_mode_tl { plain }

\tl_new:N \l_marking_role_tl
\tl_new:N \l_marking_kind_tl
\tl_new:N \l_marking_color_tl
\tl_new:N \l_marking_layout_tl
\tl_new:N \l_marking_caption_tl
\tl_new:N \l_marking_label_tl
\tl_new:N \l_marking_state_tl
\bool_new:N \l_marking_has_label_bool
\bool_new:N \g_marking_in_mark_bool

\cs_new_protected:Npn \marking_warn:nn #1#2
  { \PackageWarningNoLine { marking } { #1:~#2 } }

\cs_new:Npn \marking_normalize_token:n #1
  { \tl_lower_case:n { \tl_trim_spaces:n {#1} } }

\cs_new_protected:Npn \marking_set_fallback_color:n #1
  {
    \tl_set:Nx \l_marking_color_tl { \tl_trim_spaces:n {#1} }
    \tl_if_blank:VTF \l_marking_color_tl
      {
        \marking_warn:nn { MarkWarnInvalidColor }
          { empty~fallback~color;~keeping~current~fallback }
      }
      {
        \tl_gset_eq:NN \g_marking_fallback_color_tl \l_marking_color_tl
        % Keep LaTeX2e mirror in sync (used by biblatex citation hooks).
        \xdef\marking@fallbackcolor { \tl_use:N \l_marking_color_tl }
      }
  }

\cs_new_protected:Npn \marking_validate_state:n #1
  {
    \tl_set:Nx \l_marking_state_tl { \marking_normalize_token:n {#1} }
    \tl_if_eq:VnTF \l_marking_state_tl { colored }
      {}
      {
        \tl_if_eq:VnTF \l_marking_state_tl { plain }
          {}
          {
            \tl_if_eq:VnTF \l_marking_state_tl { hidden }
              {}
              {
                \marking_warn:nn { MarkWarnInvalidState }
                  { unknown~state~'\l_marking_state_tl'~for~role;~using~colored }
                \tl_set:Nn \l_marking_state_tl { colored }
              }
          }
      }
  }

\cs_new_protected:Npn \marking_resolve_role_state:n #1
  {
    \tl_set:Nx \l_marking_role_tl { \marking_normalize_token:n {#1} }
    \prop_get:NVNTF
      \g_marking_role_state_prop
      \l_marking_role_tl
      \l_marking_state_tl
      {}
      {
        \marking_warn:nn { MarkWarnMissingRole }
          { role~'\l_marking_role_tl'~has~no~state;~using~colored }
        \tl_set:Nn \l_marking_state_tl { colored }
      }
  }

\cs_new_protected:Npn \marking_define_role:nnn #1#2#3
  {
    \tl_set:Nx \l_marking_role_tl { \marking_normalize_token:n {#1} }
    \tl_if_blank:VTF \l_marking_role_tl
      {
        \marking_warn:nn { MarkWarnEmptyRole }
          { role~name~is~empty;~definition~ignored }
      }
      {
        \prop_if_in:NVTF \g_marking_role_color_prop \l_marking_role_tl
          {
            \marking_warn:nn { MarkWarnDuplicateRole }
              { role~'\l_marking_role_tl'~already~defined;~keeping~first~definition }
          }
          {
            \tl_set:Nx \l_marking_color_tl { \tl_trim_spaces:n {#2} }
            \tl_if_blank:VTF \l_marking_color_tl
              {
                \marking_warn:nn { MarkWarnInvalidColor }
                  { empty~color~for~role~'\l_marking_role_tl';~using~fallback }
                \prop_gput:NVV
                  \g_marking_role_color_prop
                  \l_marking_role_tl
                  \g_marking_fallback_color_tl
                % CS-name mirror: \marking_rolecolor_<role> -> color value.
                % Used by biblatex citation hooks (\marking@resolveentryroles)
                % without needing expl3 prop-list access in non-expl3 context.
                \tl_set:Nx \l_tmpa_tl
                  { marking_rolecolor_ \tl_use:N \l_marking_role_tl }
                \exp_args:Nc \xdef { \tl_use:N \l_tmpa_tl }
                  { \tl_use:N \g_marking_fallback_color_tl }
                % Store and mirror state.
                \marking_validate_state:n {#3}
                \prop_gput:NVV
                  \g_marking_role_state_prop
                  \l_marking_role_tl
                  \l_marking_state_tl
                \tl_set:Nx \l_tmpa_tl
                  { marking_rolestate_ \tl_use:N \l_marking_role_tl }
                \exp_args:Nc \xdef { \tl_use:N \l_tmpa_tl }
                  { \tl_use:N \l_marking_state_tl }
              }
              {
                \prop_gput:NVV
                  \g_marking_role_color_prop
                  \l_marking_role_tl
                  \l_marking_color_tl
                % CS-name mirror for biblatex citation hooks.
                \tl_set:Nx \l_tmpa_tl
                  { marking_rolecolor_ \tl_use:N \l_marking_role_tl }
                \exp_args:Nc \xdef { \tl_use:N \l_tmpa_tl }
                  { \tl_use:N \l_marking_color_tl }
                % Store and mirror state.
                \marking_validate_state:n {#3}
                \prop_gput:NVV
                  \g_marking_role_state_prop
                  \l_marking_role_tl
                  \l_marking_state_tl
                \tl_set:Nx \l_tmpa_tl
                  { marking_rolestate_ \tl_use:N \l_marking_role_tl }
                \exp_args:Nc \xdef { \tl_use:N \l_tmpa_tl }
                  { \tl_use:N \l_marking_state_tl }
              }
          }
      }
  }

\cs_new_protected:Npn \marking_resolve_role_color:n #1
  {
    \tl_set:Nx \l_marking_role_tl { \marking_normalize_token:n {#1} }
    \tl_if_blank:VTF \l_marking_role_tl
      {
        \marking_warn:nn { MarkWarnEmptyRole }
          { empty~role~at~call~site;~using~fallback }
        \tl_set_eq:NN \l_marking_color_tl \g_marking_fallback_color_tl
      }
      {
        \prop_get:NVNTF
          \g_marking_role_color_prop
          \l_marking_role_tl
          \l_marking_color_tl
          {}
          {
            \marking_warn:nn { MarkWarnMissingRole }
              { role~'\l_marking_role_tl'~is~undefined;~using~fallback }
            \tl_set_eq:NN \l_marking_color_tl \g_marking_fallback_color_tl
          }
      }
    \edef\MarkingResolvedColor{\tl_use:N \l_marking_color_tl}
  }

\cs_new_protected:Npn \marking_resolve_kind:n #1
  {
    \tl_set:Nx \l_marking_kind_tl { \marking_normalize_token:n {#1} }
    \tl_if_blank:VTF \l_marking_kind_tl
      { \tl_set:Nn \l_marking_kind_tl { text } }
      {}
  }

\cs_new_protected:Npn \marking_set_edit_mode:n #1
  {
    \tl_set:Nx \l_tmpa_tl { \marking_normalize_token:n {#1} }
    \str_case:VnF \l_tmpa_tl
      {
        { plain }  { \tl_gset:Nn \g_marking_edit_mode_tl { plain } }
        { popup }  { \tl_gset:Nn \g_marking_edit_mode_tl { popup } }
        { inline } { \tl_gset:Nn \g_marking_edit_mode_tl { inline } }
      }
      {
        \marking_warn:nn { MarkWarnInvalidEditMode }
          { unknown~edit~mode~'\l_tmpa_tl';~keeping~current~mode }
      }
  }

\cs_new_protected:Npn \marking_edit_plain:nn #1#2
  {
    \str_if_eq:eeTF {#1} { math } { $#2$ } {#2}
  }

\cs_new_protected:Npn \marking_edit_new:nn #1#2
  {
    \str_if_eq:eeTF {#1} { math }
      { \colorbox{\MarkingResolvedColor!25}{$\displaystyle #2$} }
      { \sethlcolor{\MarkingResolvedColor!25}\hl{#2} }
  }

\cs_new_protected:Npn \marking_edit_popup:nnn #1#2#3
  {
    \tl_if_blank:nTF {#3}
      {
        % A deletion has no replacement text to anchor a markup annotation.
        \pdfcomment[icon=Note,color=red,subject={Deleted text}]
          {Removed\space text:\textCR #2}
        \nobreak\hspace{2.2em}
      }
      {
        % PDF markup anchors are unreliable in moving caption arguments.
        \cs_if_exist:cTF {@captype}
          {
            \str_if_eq:eeTF {#1} { math }
              { \pdfcomment[icon=Note,color=red,subject={Old expression}]
                  {Earlier\space version:\textCR $#2$} }
              { \pdfcomment[icon=Note,color=red,subject={Old text}]
                  {Earlier\space version:\textCR #2} }
            \nobreak\hspace{2.2em}
            \marking_edit_new:nn {#1}{#3}
          }
          {
            \str_if_eq:eeTF {#1} { math }
              { \pdfmarkupcomment[markup=Highlight,color=\MarkingResolvedColor,
                  opacity=0.35,subject={Old expression}]
                  {$#3$}{Earlier\space version:\textCR $#2$} }
              { \pdfmarkupcomment[markup=Highlight,color=\MarkingResolvedColor,
                  opacity=0.35,subject={Old text}]
                  {#3}{Earlier\space version:\textCR #2} }
          }
      }
  }

\cs_new_protected:Npn \marking_edit_inline_old:nn #1#2
  {
    \str_if_eq:eeTF {#1} { math }
      { \textcolor{red}{\st{$#2$}} }
      { \textcolor{red}{\st{#2}} }
    \nobreak\hspace{0.2em}
  }

\cs_new_protected:Npn \marking_edit:nnnn #1#2#3#4
  {
    \marking_resolve_role_color:n {#1}
    \marking_resolve_role_state:n {#1}
    \marking_resolve_kind:n {#2}
    \tl_if_eq:VnF \l_marking_kind_tl { math }
      {
        \tl_if_eq:VnF \l_marking_kind_tl { text }
          {
            \marking_warn:nn { MarkWarnUnknownEditKind }
              { unknown~edit~kind~'\l_marking_kind_tl';~using~text }
            \tl_set:Nn \l_marking_kind_tl { text }
          }
      }
    % A plain or hidden role always leaves only the accepted replacement.
    \tl_if_eq:VnTF \l_marking_state_tl { colored }
      {
        \tl_if_eq:VnTF \g_marking_edit_mode_tl { plain }
          { \marking_edit_plain:nn {\l_marking_kind_tl}{#4} }
          {
            \tl_if_eq:VnTF \g_marking_edit_mode_tl { popup }
              {
                \tl_if_blank:nTF {#3}
                  { \marking_edit_new:nn {\l_marking_kind_tl}{#4} }
                  { \marking_edit_popup:nnn {\l_marking_kind_tl}{#3}{#4} }
              }
              {
                \tl_if_blank:nF {#3}
                  { \marking_edit_inline_old:nn {\l_marking_kind_tl}{#3} }
                \tl_if_blank:nF {#4}
                  { \marking_edit_new:nn {\l_marking_kind_tl}{#4} }
              }
          }
      }
      { \marking_edit_plain:nn {\l_marking_kind_tl}{#4} }
  }

\cs_new_protected:Npn \marking_render_inline:nnn #1#2#3
  {
    \marking_resolve_role_color:n {#1}
    \marking_resolve_role_state:n {#1}
    \marking_resolve_kind:n {#2}
    \tl_if_eq:VnTF \l_marking_kind_tl { text }
      {
        \tl_if_eq:VnTF \l_marking_state_tl { hidden }
          {}
          {
            \tl_if_eq:VnTF \l_marking_state_tl { colored }
              {\sethlcolor{\MarkingResolvedColor}\hl{#3}}
              {#3}
          }
      }
      {
        \tl_if_eq:VnTF \l_marking_kind_tl { math }
          {
            \tl_if_eq:VnTF \l_marking_state_tl { hidden }
              {}
              {
                \tl_if_eq:VnTF \l_marking_state_tl { colored }
                  { \colorbox{\MarkingResolvedColor!30}{$\displaystyle #3$} }
                  { $#3$ }
              }
          }
          {
            \tl_if_eq:VnTF \l_marking_kind_tl { figure }
              {
                \marking_warn:nn { MarkWarnInlineFigureKind }
                  { figure~kind~is~not~supported~for~Mark;~using~text }
              }
              {
                \marking_warn:nn { MarkWarnUnknownKind }
                  { unknown~kind~'\l_marking_kind_tl';~using~text }
              }
            \tl_if_eq:VnTF \l_marking_state_tl { hidden }
              {}
              {
                \tl_if_eq:VnTF \l_marking_state_tl { colored }
                  {\sethlcolor{\MarkingResolvedColor}\hl{#3}}
                  {#3}
              }
          }
      }
  }

\cs_new_protected:Npn \marking_render_heading:nn #1#2
  {
    \marking_resolve_role_state:n {#1}
    \tl_if_eq:VnTF \l_marking_state_tl { hidden }
      {
        \marking_warn:nn { MarkWarnHiddenHeading }
          { role~'\l_marking_role_tl'~is~hidden~but~used~in~a~heading;~rendering~plain }
        #2
      }
      {
        % Heading-safe rendering: highlighted in typeset text, plain in PDF strings.
        \texorpdfstring
          { \marking_render_inline:nnn {#1}{text}{#2} }
          { #2 }
      }
  }

\cs_new_protected:Npn \marking_render_block:nnn #1#2#3
  {
    \marking_resolve_role_color:n {#1}
    \marking_resolve_role_state:n {#1}
    \marking_resolve_kind:n {#2}
    \tl_if_eq:VnTF \l_marking_kind_tl { math }
      {
        \tl_if_eq:VnTF \l_marking_state_tl { hidden }
          {}
          {
            \tl_if_eq:VnTF \l_marking_state_tl { colored }
              {
                \begin{tcolorbox}[
                  enhanced,
                  colback=\MarkingResolvedColor!50,
                  colframe=\MarkingResolvedColor!50,
                  boxrule=0pt,
                  arc=0pt,
                  left=2pt,
                  right=2pt,
                  top=2pt,
                  bottom=2pt,
                  boxsep=0pt
                ]
                  \begin{equation}
                    #3
                  \end{equation}
                \end{tcolorbox}
              }
              {
                \begin{equation}
                  #3
                \end{equation}
              }
          }
      }
      {
        \tl_if_eq:VnTF \l_marking_kind_tl { figure }
          {
            \tl_if_eq:VnTF \l_marking_state_tl { hidden }
              {}
              {
                \tl_if_eq:VnTF \l_marking_state_tl { colored }
                  {
                    \begin{tcolorbox}[
                      enhanced,
                      colback=white,
                      colframe=\MarkingResolvedColor,
                      boxrule=3pt,
                      arc=0pt,
                      left=0pt,
                      right=0pt,
                      top=0pt,
                      bottom=0pt,
                      boxsep=0pt
                    ]
                      #3
                    \end{tcolorbox}
                  }
                  {#3}
              }
          }
          {
            \tl_if_eq:VnF \l_marking_kind_tl { text }
              {
                \marking_warn:nn { MarkWarnUnknownKind }
                  { unknown~kind~'\l_marking_kind_tl';~using~text }
              }
            \tl_if_eq:VnTF \l_marking_state_tl { hidden }
              {}
              {
                \tl_if_eq:VnTF \l_marking_state_tl { colored }
                  {\sethlcolor{\MarkingResolvedColor}\hl{#3}}
                  {#3}
              }
          }
      }
  }

\cs_new_protected:Npn \marking_placeholder_box:nn #1#2
  {
    \fcolorbox{\MarkingResolvedColor}{\MarkingResolvedColor!15}{
      \parbox{#1}{\centering\vspace{2em}%
        \textbf{[Placeholder Figure]}\par\medskip #2%
        \vspace{2em}}%
    }
  }

\cs_new_protected:Npn \marking_placeholder_box_plain:nn #1#2
  {
    \fbox{%
      \parbox{#1}{\centering\vspace{2em}%
        \textbf{[Placeholder Figure]}\par\medskip #2%
        \vspace{2em}}%
    }
  }

\cs_new_protected:Npn \marking_render_placeholder:nnnnn #1#2#3#4#5
  {
    \marking_resolve_role_color:n {#1}
    \marking_resolve_role_state:n {#1}
    \tl_set:Nx \l_marking_layout_tl { \marking_normalize_token:n {#2} }
    \tl_if_blank:VTF \l_marking_layout_tl
      { \tl_set:Nn \l_marking_layout_tl { single } }
      {}

    \tl_set:Nx \l_marking_caption_tl { \tl_trim_spaces:n {#5} }
    \tl_if_blank:VTF \l_marking_caption_tl
      {
        \marking_warn:nn { MarkWarnEmptyCaption }
          { placeholder~caption~is~empty;~using~'Placeholder~Figure' }
        \tl_set:Nn \l_marking_caption_tl { Placeholder Figure }
      }
      {}

    \tl_set:Nx \l_marking_label_tl { \tl_trim_spaces:n {#4} }
    \tl_if_blank:VTF \l_marking_label_tl
      {
        \marking_warn:nn { MarkWarnEmptyLabel }
          { placeholder~label~is~empty;~label~omitted }
        \bool_set_false:N \l_marking_has_label_bool
      }
      { \bool_set_true:N \l_marking_has_label_bool }

    \tl_if_eq:VnTF \l_marking_state_tl { hidden }
      {}
      {
        \tl_if_eq:VnTF \l_marking_layout_tl { wide }
          {
            \begin{figure*}[t]
              \centering
              \tl_if_eq:VnTF \l_marking_state_tl { colored }
                { \marking_placeholder_box:nn {0.85\textwidth} {#3} }
                { \marking_placeholder_box_plain:nn {0.85\textwidth} {#3} }
              \caption{\l_marking_caption_tl}
              \bool_if:NT \l_marking_has_label_bool { \exp_args:NV \label \l_marking_label_tl }
            \end{figure*}
          }
          {
            \tl_if_eq:VnF \l_marking_layout_tl { single }
              {
                \marking_warn:nn { MarkWarnUnknownLayout }
                  { unknown~layout~'\l_marking_layout_tl';~using~single }
              }
            \begin{figure}[t]
              \centering
              \tl_if_eq:VnTF \l_marking_state_tl { colored }
                { \marking_placeholder_box:nn {0.9\linewidth} {#3} }
                { \marking_placeholder_box_plain:nn {0.9\linewidth} {#3} }
              \caption{\l_marking_caption_tl}
              \bool_if:NT \l_marking_has_label_bool { \exp_args:NV \label \l_marking_label_tl }
            \end{figure}
          }
      }
  }

\cs_new_protected:Npn \marking_if_typeset:TF #1#2
  { \use:c { marking@iftypesetTF } {#1}{#2} }

\cs_new_protected:Npn \marking_mark:nnn #1#2#3
  {
    \marking_resolve_kind:n {#2}
    \tl_if_eq:VnTF \l_marking_kind_tl { heading }
      {
        \bool_if:NTF \g_marking_in_mark_bool
          {
            \marking_warn:nn { MarkWarnNestedMark }
              { nested~Mark/MarkEnv~usage;~rendering~nested~content~plain }
            #3
          }
          {
            \bool_gset_true:N \g_marking_in_mark_bool
            \marking_render_heading:nn {#1}{#3}
            \bool_gset_false:N \g_marking_in_mark_bool
          }
      }
      {
        \marking_if_typeset:TF
          {
            \bool_if:NTF \g_marking_in_mark_bool
              {
                \marking_warn:nn { MarkWarnNestedMark }
                  { nested~Mark/MarkEnv~usage;~rendering~nested~content~plain }
                #3
              }
              {
                \bool_gset_true:N \g_marking_in_mark_bool
                \marking_render_inline:nnn {#1}{#2}{#3}
                \bool_gset_false:N \g_marking_in_mark_bool
              }
          }
          {
            \marking_warn:nn { MarkWarnMovingArgFallback }
              { Mark~used~in~moving~argument;~using~plain~text }
            #3
          }
      }
  }

\cs_new_protected:Npn \marking_markenv:nnn #1#2#3
  {
    \bool_if:NTF \g_marking_in_mark_bool
      {
        \marking_warn:nn { MarkWarnNestedMark }
          { nested~Mark/MarkEnv~usage;~rendering~nested~content~plain }
        #3
      }
      {
        \bool_gset_true:N \g_marking_in_mark_bool
        \marking_render_block:nnn {#1}{#2}{#3}
        \bool_gset_false:N \g_marking_in_mark_bool
      }
  }

\ExplSyntaxOff

\ExplSyntaxOn
\NewDocumentCommand{\SetMarkFallbackColor}{m}
  { \marking_set_fallback_color:n {#1} }

\NewDocumentCommand{\DefineMarkRole}{m m O{colored}}
  { \marking_define_role:nnn {#1}{#2}{#3} }

\NewDocumentCommand{\Mark}{m O{text} +m}
  { \marking_mark:nnn {#1}{#2}{#3} }

\NewDocumentCommand{\SetMarkEditMode}{m}
  { \marking_set_edit_mode:n {#1} }

\NewDocumentCommand{\MarkEdit}{m O{text} +m +m}
  { \marking_edit:nnnn {#1}{#2}{#3}{#4} }

\NewDocumentEnvironment{MarkEnv}{m O{text} +b}
  { \marking_markenv:nnn {#1}{#2}{#3} }
  {}

\NewDocumentCommand{\MarkPlaceholderFigure}{m O{single} m m m}
  { \marking_render_placeholder:nnnnn {#1}{#2}{#3}{#4}{#5} }
\ExplSyntaxOff

\makeatletter
\AtBeginDocument{%
  \@ifpackageloaded{hyperref}{%
    \pdfstringdefDisableCommands{%
      % In PDF strings/bookmarks, keep only the marked content.
      % Heading marks always include [heading] in this project.
      \def\Mark#1#2]#3{#3}%
    }%
  }{}%
}
\makeatother

% ============================================================================
% Citation Marking (biblatex-only feature)
% Soft deps: biblatex, hyperref, xpatch — degrade gracefully when absent.
% Usage: add  markroles = {role}  to a .bib entry to highlight it.
%        Role names must match a \DefineMarkRole definition.
% ============================================================================

% breakable is needed for page-spanning bib panels; call is idempotent.
\tcbuselibrary{breakable}

\makeatletter
% LaTeX2e mirror of \g_marking_fallback_color_tl; kept in sync by
% \marking_set_fallback_color:n (via \xdef). Used by biblatex citation hooks.
\newcommand{\marking@fallbackcolor}{red}
% markroles data model registration lives in markings.dbx (loaded via
% biblatex's datamodel= package option). See publications/modules/markings/markings.dbx.

% \SetMarkDefaultCiteColor{<color>}
% Optional public override for the reset citecolor used in \AtEveryCitekey.
% Normally auto-captured from hyperref at \AtBeginDocument; only set this if
% your project changes citecolor after that point.
\newcommand{\marking@defaultcitecolor}{} % placeholder; filled at \AtBeginDocument
\newcommand{\SetMarkDefaultCiteColor}[1]{\renewcommand{\marking@defaultcitecolor}{#1}}

\AtBeginDocument{%
  \@ifpackageloaded{biblatex}{%
    %
    % ---- Helper: resolve markroles field and set \marking@entryrolecolor ----
    % Sets \marking@entryrolecolor to the resolved color, or \@empty if none.
    % markroles is declared as a scalar literal field (markings.dbx), accessed
    % via \thefield.  The role-to-color lookup uses a CS-name mirror defined
    % by \DefineMarkRole:  \csname marking_rolecolor_<role>\endcsname.
    % The role-to-state lookup uses \csname marking_rolestate_<role>\endcsname.
    % Both mirrors avoid expl3 prop-list access from non-expl3 macro bodies.
    % \marking@statecolored is a constant used for \ifx comparisons.
    \def\marking@statecolored{colored}%
    \newcommand{\marking@resolveentryroles}{%
      \global\let\marking@entryrolecolor\@empty
      \global\def\marking@entryrolestate{colored}%
      \iffieldundef{markroles}{}{%
        \edef\marking@tmprole{\thefield{markroles}}%
        \@ifundefined{marking_rolecolor_\marking@tmprole}{%
          \PackageWarningNoLine{marking}{MarkWarnMissingRole:
            role '\marking@tmprole' is undefined; using fallback}%
          \global\edef\marking@entryrolecolor{\marking@fallbackcolor}%
        }{%
          \global\edef\marking@entryrolecolor{%
            \csname marking_rolecolor_\marking@tmprole\endcsname}%
        }%
        \@ifundefined{marking_rolestate_\marking@tmprole}{}{%
          \global\edef\marking@entryrolestate{%
            \csname marking_rolestate_\marking@tmprole\endcsname}%
        }%
      }%
    }%
    %
    % ---- Hook 1: in-text citation color (requires hyperref) ----
    \@ifpackageloaded{hyperref}{%
      % Auto-capture the project's default citecolor from hyperref internals.
      \@ifundefined{HyColor@citecolor}{%
        \@ifundefined{Hy@citecolor}{%
          \renewcommand{\marking@defaultcitecolor}{black}%
        }{%
          \edef\marking@defaultcitecolor{\Hy@citecolor}%
        }%
      }{%
        \edef\marking@defaultcitecolor{\HyColor@citecolor}%
      }%
      %
      \AtEveryCitekey{%
        \marking@resolveentryroles
        \ifx\marking@entryrolecolor\@empty
          \hypersetup{citecolor=\marking@defaultcitecolor}%
        \else
          \ifx\marking@entryrolestate\marking@statecolored
            \hypersetup{citecolor=\marking@entryrolecolor!50!black}%
          \else
            \hypersetup{citecolor=\marking@defaultcitecolor}%
          \fi
        \fi
      }%
    }{%
      \PackageInfo{marking}{hyperref not loaded; in-text citation color disabled}%
    }%
    %
    % ---- Hook 2: bibliography label box ----
    \DeclareFieldFormat{labelnumberwidth}{%
      \marking@resolveentryroles
      \ifx\marking@entryrolecolor\@empty
        \mkbibbrackets{#1}%
      \else
        \ifx\marking@entryrolestate\marking@statecolored
          \colorbox{\marking@entryrolecolor!25}{\strut\mkbibbrackets{#1}}%
        \else
          \mkbibbrackets{#1}%
        \fi
      \fi
    }%
    %
    % ---- Hook 3: bibliography entry body panel (requires xpatch) ----
    \@ifpackageloaded{xpatch}{%
      \xpretobibmacro{begentry}{%
        \marking@resolveentryroles
        \ifx\marking@entryrolecolor\@empty\else
          \ifx\marking@entryrolestate\marking@statecolored
            \begin{tcolorbox}[%
              enhanced, breakable,%
              colback=\marking@entryrolecolor!25,%
              colframe=\marking@entryrolecolor,%
              boxrule=1pt, arc=0pt,%
              left=0pt, right=0pt, top=0pt, bottom=0pt, boxsep=1pt,%
              before skip=0pt, after skip=0pt,%
              enlarge left by=-\labelsep%
            ]\ignorespaces
          \fi
        \fi
      }{}{}%
      \xapptobibmacro{finentry}{%
        \marking@resolveentryroles
        \ifx\marking@entryrolecolor\@empty\else
          \ifx\marking@entryrolestate\marking@statecolored
            \unskip\end{tcolorbox}%
          \fi
        \fi
      }{}{}%
    }{%
      \PackageInfo{marking}{xpatch not loaded; bib entry body panels disabled}%
    }%
  }{%
    \PackageInfo{marking}{biblatex not loaded; citation marking disabled}%
  }%
}
\makeatother


% Role colors. Highlights use the role color at full strength, so keep it light.
\definecolor{MarkChanges}{RGB}{150, 220, 150}
\definecolor{MarkMissing}{RGB}{245, 130, 130}

% Review build: every role is colored.
\DefineMarkRole{changes}{MarkChanges}
\DefineMarkRole{edits}{MarkChanges}
\DefineMarkRole{missing}{MarkMissing}

% How \MarkEdit shows replacements: plain (new text only, the module default),
% popup (new text highlighted, old text in a PDF note), or inline (old text
% struck through in red, new text highlighted).
\SetMarkEditMode{popup}

% Submission build: replace the three role lines above with these, so marked
% text prints without highlights and replacements show only the new text.
% \DefineMarkRole{changes}{MarkChanges}[plain]
% \DefineMarkRole{edits}{MarkChanges}[plain]
% \DefineMarkRole{missing}{MarkMissing}[plain]

% ---------- Math operators ----------
\DeclareMathOperator*{\argmin}{arg\,min}
\DeclareMathOperator{\E}{\mathbb{E}}
\DeclareMathOperator{\Var}{Var}

% ---------- Theorem environments ----------
\newtheorem{definition}{Definition}
\newtheorem{proposition}{Proposition}
\newtheorem{remark}{Remark}

% ---------- Author block ----------
\usepackage{authblk}
